In questo capitolo si completa l'analisi prosodica analizzando il testo dal punto di vista emotivo e semantico. 
Questa fase è particolarmente importante per la successiva  generazione musicale che guida e orienta l'atmosfera melodica e la scelta automatica 
degli strumenti. A tale scopo, il sistema elabora il testo di input utilizzando un approccio ibrido che fonde risorse lessicali interne al programma 
con database esterni. Per rappresentare le emozioni è stato predisposto un modello dimensionale, che mappa valori numerici estrapolati dalle singole 
parole in uno spazio vettoriale continuo (con valori normalizzati tra -1.0 e +1.0) lungo tre assi fondamentali: 
\begin{itemize}
    \item \textbf{Valenza} (polarità positiva/negativa): estendendosi da stati di dolore a stati di gioia, questo asse guida principalmente le scelte modali 
    e armoniche (es. l'utilizzo di scale maggiori o minori).
    \item \textbf{Arousal} (energia/calma): misura il livello di attivazione e tensione. In ambito musicale, si traduce tipicamente in variazioni di tempo, 
    densità ritmica e intensità dinamica.
    \item \textbf{Tenerezza} (affetto/distacco): asse cruciale per il dominio poetico, permette di isolare sentimenti di dolcezza e intimità, distinguendo 
    ad esempio un brano genericamente positivo da uno profondamente affettuoso, influenzando la scelta dei timbri strumentali.
\end{itemize}



\section{Risorse lessicali}
Il calcolo dell'\textit{embedding emotivo} si fonda su un approccio \textit{dictionary-based}, in cui ogni termine del testo viene cercato all'interno di 
specifici lessici per ottenerne le coordinate dimensionali. I dizionari curati manualmente offrono altissima precisione per un dominio specifico, ma hanno 
una copertura limitata; al contrario, i grandi lessici del Natural Language Processing (NLP) offrono una copertura vastissima, ma possono risultare 
imprecisi o sbilanciati rispetto alle necessità del dominio poetico-musicale. Per risolvere questo \textit{trade-off}, il sistema implementa un'architettura 
ibrida che interroga gerarchicamente tre diverse fonti di dati: un dizionario interno dedicato e due risorse esterne appartenenti al progetto NRC 
(National Research Council Canada).

\subsection{Dizionario interno (Lexicon)}
Il nucleo dell'analisi è costituito dal dizionario interno (\texttt{LEXICON}), una struttura dati costruita e tarata ad hoc per il software. 
Ogni lemma al suo interno è associato a una tupla di tre valori compresi tra -1 e +1, corrispondenti rispettivamente a Valenza, Arousal e Tenerezza. 
Questo modulo è fornito di supporto a quelli esterni estremamente più completi e precisi ma poveri di parole in sfumature nei quadranti a valenza positiva.
Per ovviare questo sbilanciamento e permettere al motore di generazione musicale di esplorare l'intero spettro sonoro, il dizionario interno è stato 
integrato con lemmi mirati a coprire specifiche aree del piano Valenza/Arousal: 
\begin{itemize}
    \item Calma e Serenità (V+, A-): termini come \textit{calmo, pacifico, silenzioso, adagio}, fondamentali per generare atmosfere musicali distese.
    \item Gioia ad alta intensità (V+, A+): termini come \textit{entusiasmo, trionfo o danza}, per stimolare generazioni musicali energiche e brillanti.
    \item  Affetto (Tenerezza): l'inserimento manuale del terzo asse per lemmi come \textit{amore, carezza o sogno} garantisce che il calore del testo venga 
    tradotto correttamente in musica, distinguendolo da parole genericamente positive ma fredde.
\end{itemize}

\subsection{Modello continuo: NRC VAD Lexicon}
Qualora un termine non sia presente nel dizionario interno (dopo averne estratto il lemma tramite la libreria \texttt{simplemma}), il sistema interroga il suo 
fallback primario: l'NRC VAD Lexicon.\\
A differenza di altri dizionari basati su categorie discrete, il VAD Lexicon contiene punteggi continui annotati direttamente da esseri umani per i tre 
assi di Valenza, Arousal e Dominanza. Questa caratteristica lo rende strutturalmente affine al modello dimensionale implementato.\\
L'uso di misurazioni dirette risolve nativamente il bias dell'Arousal tipico di altri database: non dovendo derivare l'intensità energetica da 
etichette predefinite, il VAD Lexicon mappa con estrema precisione anche le parole a bassissima energia (come calma, sonno, quiete).\\
A livello implementativo, il modulo relativo è stato dotato di una funzione di rilevamento automatico della scala metrica. Poiché le diverse versioni 
storiche del dataset NRC (2018 vs 2025) utilizzano scale differenti (da [0, 1] a [-1, 1]), il loader normalizza automaticamente i valori sulla 
scala standard del software (da -1 a +1) applicando la trasformazione lineare $v_{\mathit{new}} = v_{\mathit{old}} \times 2 - 1$, garantendo l'omogeneità 
dei dati in ingresso.

\subsection{Modello categoriale: NRC Emotion Lexicon (EmoLex)}       
Qualora il lemma risulti assente anche nel VAD, il sistema ricorre al fallback secondario: l'NRC Emotion Lexicon. Questo database è inoltre impiegato 
come sorgente esterna esclusiva per derivare l'asse della Tenerezza, non presente nel dataset VAD.\\
L'integrazione dell'EmoLex ha richiesto un importante lavoro di adattamento algoritmico. Essendo un lessico categoriale basato sulle 8 emozioni di base 
di Plutchik, le etichette binarie (es. \textit{joy=1}, \textit{anger=1}) sono state matematicamente proiettate sullo spazio vettoriale di Russell 
assegnando coordinate costanti a ciascuna categoria (\texttt{CATEGORY\_TO\_VA}). Tuttavia, l'uso dell'EmoLex nella sua traduzione italiana presentava 
sbilanciamenti sistemici (il risultato tendeva a posizionarsi sempre su un Arousal positivo). Per ovviare a questo problema, il modulo implementa 
due euristiche correttive:
\begin{itemize}
    \item Ponderazione dell'affidabilità categoriale (\texttt{CATEGORY\_RELIABILITY}): L'EmoLex italiano, essendo generato per traduzione automatica dal lemma inglese, 
    presenta "rumore" semantico. L'algoritmo non effettua una media semplice delle categorie attive, ma applica pesi ridotti a quelle più rumorose, 
    permettendo a emozioni più stabili e ad Arousal negativo (come sadness) di prevalere nel calcolo vettoriale.
    \item Correzione della polarità grezza (\texttt{POLARITY\_TO\_VA}): Quando una parola presenta solo un flag grezzo negative (senza emozioni specifiche), 
    l'Arousal di default è stato forzato in negativo (-0.15). Questa scelta riflette l'assunto psico-linguistico per cui la negatività "pura" 
    si avvicini maggiormente a stati di letargia o tristezza, piuttosto che all'alta attivazione della rabbia vincolati a  valori di Arousal positivi.
\end{itemize}

\subsection{Logica di fusione e assegnazione dei pesi}
Infine viene gestita l'estrazione vettoriale finale per una singola parola. Per massimizzare le probabilità di match, il sistema tenta prima di cercare 
il lemma della parola e, in caso di fallimento, tenta con la parola flessa originale.\\
La criticità principale si verifica quando un termine è presente contemporaneamente nel dizionario interno e in una delle risorse esterne (VAD o EmoLex). 
Sostituire del tutto il valore generico esterno al valore manuale vanificherebbe il lavoro di calibrazione fatto sul dominio poetico. Pertanto, l'algoritmo 
fonde i dati applicando una media pesata lineare. I pesi sono stati definiti empiricamente nel codice:
\begin{itemize}
    \item \texttt{HAND\_WEIGHT = 0.6}
    \item \texttt{NRC\_WEIGHT = 0.4}
\end{itemize}
Questa proporzione (60\% al lessico manuale, 40\% a quello NLP) garantisce che il dizionario interno LEXICON mantenga la priorità assoluta, 
fungendo da "centro di gravità" per l'interpretazione emotiva della parola, pur permettendo al dato esterno di smussare il risultato finale, 
arricchendolo con le sfumature statistiche calcolate su vasta scala.

\section{Pipeline di analisi emotiva}
Una volta stabilite e ponderate le risorse lessicali, il sistema elabora l'intero testo in ingresso per calcolarne l'emotività globale. 
Questo processo non consiste in una semplice media matematica delle parole trovate, ma richiede una pipeline algoritmica divisa in fasi logiche sequenziali.
\begin{enumerate}
    \item \textbf{Preprocessing e normalizzazione linguistica}: il testo viene convertito in minuscolo, privato della punteggiatura e suddiviso in 
    token isolati. Successivamente, ogni parola viene sottoposta a lemmatizzazione tramite la libreria \texttt{simplemma}. Questo passaggio è fondamentale 
    per il trattamento della lingua italiana: riconducendo i verbi coniugati all'infinito e le parole flesse al singolare maschile, si massimizza 
    la probabilità di riscontro all'interno dei dizionari vettoriali. Contestualmente, i termini puramente funzionali privi di carico emotivo 
    (come articoli, preposizioni, congiunzioni e forme verbali comuni) vengono scartati tramite un set di \texttt{STOPWORDS}, affinché non alterino i 
    calcoli successivi.
    \item \textbf{Valutazione del contesto sintattico} (modificatori): l'analisi semantica non si limita al mero rintracciamento lessicale (approccio \textit{bag of words}): 
    il significato emotivo di una parola dipende fortemente dai termini che la precedono. Il modulo gestisce questa dinamica in modo stateful, 
    avvalendosi di costrutti logici attivati da specifici modificatori:
    \begin{itemize}
        \item Negazioni: quando l'algoritmo incontra termini come "non" o "senza", attiva un flag che modifica i valori del bersaglio successivo. 
        Valenza e Arousal vengono invertiti di segno e attenuati (rispettivamente del 60\% e 50\%), poiché negare una parola come "felice" non equivale a
        un odio assoluto, ma genera una negatività moderata. La Tenerezza subisce invece una forte smorzatura verso lo zero, indicando distacco o assenza di 
        calore.
        \item intensificatori: L'incontro con termini come "molto" o "profondamente" incrementa del 30\% le coordinate vettoriali della parola successiva. 
        Per garantire robustezza all'analisi, il parser riconosce un'eccezione definita \texttt{SCOPE\_NEUTRAL}: avverbi comuni che si interpongono tra modificatore 
        e bersaglio (come "veramente" o "davvero") vengono ignorati, permettendo allo scope del modificatore di persistere intatto. 
    \end{itemize}
    \item \textbf{Agregazione vettoriale basata su magnitude}: l'ultimo passaggio fonde i vettori delle singole parole in un unico embedding emotivo 
    rappresentativo dell'intera frase. L'uso di una media aritmetica standard comporterebbe un grave difetto strutturale: una parola emotivamente 
    fortissima verrebbe drasticamente diluita se accompagnata da termini neutri. Per ovviare a ciò, l'algoritmo calcola una media pesata, 
    dove il peso assegnato a ogni parola è proporzionale al valore assoluto più alto raggiunto sui tre assi, assicurando che i termini emotivamente 
    più intensi fungano da traino per il risultato. Per impedire errori matematici o annacquamenti, vengono applicate una tolleranza (epsilon), 
    che ignora i contributi vettoriali troppo vicini allo zero, e un peso minimo (floor), che impedisce l'azzeramento del contributo per termini a 
    carica modesta.
\end{enumerate}
L'output finale della pipeline restituisce un dizionario contenente le tre coordinate aggregate (Valenza, Arousal, Tenerezza), le parole matchate 
e un indicatore quantitativo di coverage, calcolato come rapporto tra le parole trovate nei dizionari e il totale dei token di contenuto analizzati.

\section{Rappresentazione Visiva e Proiezione Spaziale}
Al termine della pipeline di calcolo, l'algoritmo produce un \textit{embedding emotivo} numerico. Sebbene questi dati (Valenza, Arousal, Tenerezza) 
siano progettati per essere passati direttamente al motore di generazione musicale, il sistema integra un modulo di plotting basato sulla libreria 
\texttt{matplotlib}, utile a fornire un riscontro grafico e intuitivo del posizionamento emotivo del testo elaborato.\\
I risultati quantitativi dell'analisi vengono tradotti in un diagramma a dispersione (\textit{scatter plot}) bidimensionale, modellato fedelmente sul 
Modello Circomplesso di Russell. L'ambiente grafico è strutturato su un piano cartesiano limitato nell'intervallo $[-1.1, +1.1]$, in cui:
\begin{itemize}
    \item \textbf{L'asse delle ascisse (X)} rappresenta la \textbf{Valenza}, estendendosi da un polo di massima negatività (sinistra) a uno di massima 
    positività (destra).
    \item \textbf{L'asse delle ordinate (Y)} rappresenta l'\textbf{Arousal}, estendendosi da stati di letargia/calma (basso) a stati di estrema 
    tensione/intensità (alto). 
\end{itemize}

Gli assi ortogonali dividono il piano in quattro quadranti semantici ben definiti, che il software etichetta in trasparenza per facilitare la lettura 
dei risultati:
\begin{enumerate}
    \item \textbf{V+, A+ (Primo quadrante):} \textit{Gioia / Eccitazione} 
    \item \textbf{V-, A+ (Secondo quadrante):} \textit{Rabbia / Tensione}
    \item \textbf{V-, A- (Terzo quadrante):} \textit{Tristezza / Malinconia}
    \item \textbf{V+, A- (Quarto quadrante):} \textit{Calma / Serenità}
\end{enumerate}

L'intera poesia o testo in ingresso viene aggregato e tracciato sul piano sotto forma di un nodo (evidenziato in rosso), a cui vengono accostate le 
coordinate esatte.\\
È importante sottolineare una specifica scelta di design nella visualizzazione: sebbene l'algoritmo di estrazione elabori l'emozione su tre assi, 
il grafico opera una riduzione dimensionale proiettando il risultato esclusivamente in 2D (Valenza e Arousal). L'asse della Tenerezza, pur influenzando 
internamente la dinamica e il timbro della generazione sonora, viene escluso dalla rappresentazione spaziale per preservare la leggibilità del diagramma 
ed evitare la complessità visiva dei grafici tridimensionali.\\
Infine, il modulo visivo non si limita a mappare l'emozione testuale, ma anticipa il collegamento con il dominio musicale. All'interno della legenda del 
grafico viene infatti riportata dinamicamente la scala o \textit{modalità musicale} (es. Maggiore, Minore) che il sistema ha derivato dalle coordinate vettoriali, 
rendendo esplicito il ponte logico tra l'analisi semantica e le conseguenti scelte compositive.

\begin{figure}[htbp]
    \centering
    % Inserisci qui il percorso o il nome della tua immagine
    \includegraphics[width=0.75\textwidth]{Figure/emotion_space.png}
    
    \caption{Esempio di output generato dal modulo \texttt{visualizer.py}. Il grafico mostra il posizionamento vettoriale della poesia analizzata e la relativa scala musicale assegnata.}
    \label{fig:emotion_space}
\end{figure}

